\section{Circuitos y conexiones}

\subsection{ATmega328P Maestro}
\begin{table}[H]
\centering
\small
\begin{tabular}{lll}
\toprule
Pin Nano & Pin ATmega328P & Función \\
\midrule
D0 & PD0 & RX UART \\
D1 & PD1 & TX UART \\
D2--D7 & PD2--PD7 & LCD D0--D5 \\
D8 & PB0 & LCD D6 \\
D9 & PB1 & LCD D7 \\
D10 & PB2 & XSHUT VL53L0X \\
A0 & PC0 & Enable LCD \\
A1 & PC1 & RS LCD \\
A4 & PC4 & SDA I2C \\
A5 & PC5 & SCL I2C \\
\bottomrule
\end{tabular}
\caption{Asignación principal de pines del Maestro.}
\end{table}

\subsection{Slave 1 -- 0x11}
\begin{table}[H]
\centering
\small
\begin{tabular}{lll}
\toprule
Pin Nano & Pin ATmega328P & Función \\
\midrule
A2 & PC2 & ECHO HC-SR04 \\
A3 & PC3 & TRIG HC-SR04 \\
A4 & PC4 & SDA I2C \\
A5 & PC5 & SCL I2C \\
D11 & PB3 & Servo de agua \\
D9 & PB1 & ENA L298N / PWM \\
D7 & PD7 & IN1 L298N \\
D8 & PB0 & IN2 L298N \\
\bottomrule
\end{tabular}
\caption{Asignación principal de pines del Slave 1.}
\end{table}

\subsection{Slave 2 -- 0x12}
\begin{table}[H]
\centering
\small
\begin{tabular}{lll}
\toprule
Pin Nano & Pin ATmega328P & Función \\
\midrule
A0 & PC0 & Sensor infrarrojo \\
A4 & PC4 & SDA I2C \\
A5 & PC5 & SCL I2C \\
D2 & PD2 & ULN2003 IN1 \\
D3 & PD3 & ULN2003 IN2 \\
D4 & PD4 & ULN2003 IN3 \\
D5 & PD5 & ULN2003 IN4 \\
D6 & PD6 & Servo de puerta \\
\bottomrule
\end{tabular}
\caption{Asignación principal de pines del Slave 2.}
\end{table}

El bus I2C comparte SDA y SCL entre los tres ATmega328P y el VL53L0X. Se conectan resistencias de 470~\(\Omega\) en configuración pull-up en las líneas SDA y SCL. El sistema debe mantener una referencia común de tierra. Los servos, el motor DC y el Stepper requieren una distribución de potencia apropiada para evitar que las corrientes de los actuadores afecten la lógica o la comunicación I2C. La interfaz con el ESP32 considera además la diferencia entre los niveles lógicos de 5 V y 3.3 V. El TX del Maestro (D1) se dirige al RX2 del ESP32 (GPIO16), mientras que el TX2 del ESP32 (GPIO17) se dirige al receptor por software del Maestro en D12. La señal de 5 V hacia el ESP32 debe pasar por la adaptación de nivel correspondiente.


\subsection{Conexión del ESP32}
\begin{table}[H]
\centering
\caption{Señales utilizadas entre el Maestro y el ESP32.}
\begin{tabular}{lll}
\toprule
Origen & Destino & Función \\
\midrule
Master D1 / PD1 / TX & ESP32 GPIO16 / RX2 & Telemetría hacia Adafruit IO \\
ESP32 GPIO17 / TX2 & Master D12 / PB4 & Órdenes desde Adafruit IO \\
GND Master & GND ESP32 & Referencia común \\
\bottomrule
\end{tabular}
\end{table}
